\documentclass[10pt,a4paper]{amsart}
\usepackage{amsmath,amssymb,amsthm,graphicx}
\usepackage{verbatim}
\newenvironment{reference}{\comment}{\endcomment}
%\newenvironment{reference}{}{}
\newenvironment{slogan}{\comment}{\endcomment}
\newenvironment{history}{\comment}{\endcomment}

% For commutative diagrams we use Xy-pic
\usepackage[all]{xy}

% We use 2cell for 2-commutative diagrams.
\xyoption{2cell}
% Theorem environments.
%
\theoremstyle{plain}
\newtheorem{theorem}[subsection]{Theorem}
\newtheorem{proposition}[subsection]{Proposition}
\newtheorem{lemma}[subsection]{Lemma}

\theoremstyle{definition}
\newtheorem{definition}[subsection]{Definition}
\newtheorem{example}[subsection]{Example}
\newtheorem{exercise}[subsection]{Exercise}
\newtheorem{situation}[subsection]{Situation}

\theoremstyle{remark}
\newtheorem{remark}[subsection]{Remark}
\newtheorem{remarks}[subsection]{Remarks}

\numberwithin{equation}{subsection}

% Macros
%
\def\lim{\mathop{\mathrm{lim}}\nolimits}
\def\colim{\mathop{\mathrm{colim}}\nolimits}
\def\Spec{\mathop{\mathrm{Spec}}}
\def\Hom{\mathop{\mathrm{Hom}}\nolimits}
\def\Ext{\mathop{\mathrm{Ext}}\nolimits}
\def\SheafHom{\mathop{\mathcal{H}\!\mathit{om}}\nolimits}
\def\SheafExt{\mathop{\mathcal{E}\!\mathit{xt}}\nolimits}
\def\Sch{\mathit{Sch}}
\def\Mor{\mathop{\mathrm{Mor}}\nolimits}
\def\Ob{\mathop{\mathrm{Ob}}\nolimits}
\def\Sh{\mathop{\mathit{Sh}}\nolimits}
\def\NL{\mathop{N\!L}\nolimits}
\def\CH{\mathop{\mathrm{CH}}\nolimits}
\def\proetale{{pro\text{-}\acute{e}tale}}
\def\etale{{\acute{e}tale}}
\def\QCoh{\mathit{QCoh}}
\def\Ker{\mathop{\mathrm{Ker}}}
\def\Im{\mathop{\mathrm{Im}}}
\def\Coker{\mathop{\mathrm{Coker}}}
\def\Coim{\mathop{\mathrm{Coim}}}

% Boxtimes
%
\DeclareMathSymbol{\boxtimes}{\mathbin}{AMSa}{"02}

%
% Macros for moduli stacks/spaces
%
\def\QCohstack{\mathcal{QC}\!\mathit{oh}}
\def\Cohstack{\mathcal{C}\!\mathit{oh}}
\def\Spacesstack{\mathcal{S}\!\mathit{paces}}
\def\Quotfunctor{\mathrm{Quot}}
\def\Hilbfunctor{\mathrm{Hilb}}
\def\Curvesstack{\mathcal{C}\!\mathit{urves}}
\def\Polarizedstack{\mathcal{P}\!\mathit{olarized}}
\def\Complexesstack{\mathcal{C}\!\mathit{omplexes}}
% \Pic is the operator that assigns to X its picard group, usage \Pic(X)
% \Picardstack_{X/B} denotes the Picard stack of X over B
% \Picardfunctor_{X/B} denotes the Picard functor of X over B
\def\Pic{\mathop{\mathrm{Pic}}\nolimits}
\def\Picardstack{\mathcal{P}\!\mathit{ic}}
\def\Picardfunctor{\mathrm{Pic}}
\def\Deformationcategory{\mathcal{D}\!\mathit{ef}}
\usepackage{cleveref}
\crefname{originallemma}{original lemma}{original lemmas}
\Crefname{originallemma}{Original lemma}{Original lemmas}
\makeatletter
\newlabel{cohomology-lemma-Rhom-complex-of-direct-summands-finite-free}{{08I5}{source}{}{sourcepointer-1}{}}
\newlabel{cohomology-lemma-Rhom-complex-of-direct-summands-finite-free@cref}{{[originallemma][0][]08I5}{[1][0][]source}{}{}{}}
\makeatother
\crefname{originallemma}{original lemma}{original lemmas}
\Crefname{originallemma}{Original lemma}{Original lemmas}
\makeatletter
\newlabel{cohomology-lemma-Rhom-strictly-perfect}{{08DM}{source}{}{sourcepointer-2}{}}
\newlabel{cohomology-lemma-Rhom-strictly-perfect@cref}{{[originallemma][0][]08DM}{[1][0][]source}{}{}{}}
\makeatother
\crefname{originalequation}{original equation}{original equations}
\Crefname{originalequation}{Original equation}{Original equations}
\makeatletter
\newlabel{equation-epsilon}{{08H7}{source}{}{sourcepointer-3}{}}
\newlabel{equation-epsilon@cref}{{[originalequation][0][]08H7}{[1][0][]source}{}{}{}}
\makeatother
\crefname{originalreference}{original reference}{original references}
\Crefname{originalreference}{Original reference}{Original references}
\makeatletter
\newlabel{lemma-epsilon-flat}{{08H8}{source}{}{sourcepointer-4}{}}
\newlabel{lemma-epsilon-flat@cref}{{[originalreference][0][]08H8}{[1][0][]source}{}{}{}}
\makeatother
\crefname{originallemma}{original lemma}{original lemmas}
\Crefname{originallemma}{Original lemma}{Original lemmas}
\makeatletter
\newlabel{perfect-lemma-affine-compare-bounded}{{06Z0}{source}{}{sourcepointer-5}{}}
\newlabel{perfect-lemma-affine-compare-bounded@cref}{{[originallemma][0][]06Z0}{[1][0][]source}{}{}{}}
\makeatother
\crefname{originallemma}{original lemma}{original lemmas}
\Crefname{originallemma}{Original lemma}{Original lemmas}
\makeatletter
\newlabel{perfect-lemma-lift-pseudo-coherent}{{08EE}{source}{}{sourcepointer-6}{}}
\newlabel{perfect-lemma-lift-pseudo-coherent@cref}{{[originallemma][0][]08EE}{[1][0][]source}{}{}{}}
\makeatother
\crefname{originallemma}{original lemma}{original lemmas}
\Crefname{originallemma}{Original lemma}{Original lemmas}
\makeatletter
\newlabel{perfect-lemma-perfect-affine}{{08EB}{source}{}{sourcepointer-7}{}}
\newlabel{perfect-lemma-perfect-affine@cref}{{[originallemma][0][]08EB}{[1][0][]source}{}{}{}}
\makeatother
\crefname{originallemma}{original lemma}{original lemmas}
\Crefname{originallemma}{Original lemma}{Original lemmas}
\makeatletter
\newlabel{sites-cohomology-lemma-Rhom-complex-of-direct-summands-finite-free}{{08JI}{source}{}{sourcepointer-8}{}}
\newlabel{sites-cohomology-lemma-Rhom-complex-of-direct-summands-finite-free@cref}{{[originallemma][0][]08JI}{[1][0][]source}{}{}{}}
\makeatother
\crefname{originalremark}{original remark}{original remarks}
\Crefname{originalremark}{Original remark}{Original remarks}
\makeatletter
\newlabel{sites-cohomology-remark-prepare-fancy-base-change}{{08JF}{source}{}{sourcepointer-9}{}}
\newlabel{sites-cohomology-remark-prepare-fancy-base-change@cref}{{[originalremark][0][]08JF}{[1][0][]source}{}{}{}}
\makeatother
\title{Derived Hom commutes with the localization functor for pseudo-coherent or perfect first argument}
\author{The Stacks Project Authors}
\date{stacks-project:spaces-perfect.tex:local-snapshot-20261012; GNU FDL 1.2 or later}
\begin{document}
\maketitle
\noindent\textit{Editorial note.} All delivered mathematical text is the author's own, copied from the pinned Stacks Project chapter file spaces-perfect.tex; no statement or proof step was rewritten, reordered or removed.\par\medskip
\noindent\textit{Editorial note.} The counted proof is the author's complete argument: both cases are proved by explicit complex representations and reduce to the same finite-free pointwise isomorphism; the argument is complete\par\medskip
\noindent\textit{Editorial note.} The label lemma-descend-RHom names the counted statement, and the proof environment immediately following it in the source is delivered verbatim as its proof.\par\medskip
\noindent\textit{Editorial note.} The counted statement is carried by the author's own lemma environment in the chapter source.\par\medskip
\noindent\textit{Allowed standard prerequisites (editorial).}\begin{itemize}
\item the localization functor epsilon\textasciicircum{}* from O\_X-modules to O\_etale-modules on the small etale site and its flatness
\item RSheafHom in the derived category of modules on a ringed site
\item pseudo-coherent and perfect complexes represented by finite free resp. strictly perfect complexes
\item the Rhom representation lemmas for complexes of direct summands of finite free modules and for strictly perfect complexes
\end{itemize}\medskip
\ifdefined\hypertarget\hypertarget{sourcepointer-1}{}\fi
\noindent\textit{Source reference.} Original lemma 08I5: sources/stacks-chow/tags.txt, source lines 10803--10803. Prerequisite pointer invoked at proof line(s) 3456; it names a standard fact of a sibling chapter that the counted argument uses, and every step of the counted proof built on it is retained verbatim in the delivered body.\par\medskip
\ifdefined\hypertarget\hypertarget{sourcepointer-2}{}\fi
\noindent\textit{Source reference.} Original lemma 08DM: sources/stacks-chow/tags.txt, source lines 10642--10642. Prerequisite pointer invoked at proof line(s) 3495; it names a standard fact of a sibling chapter that the counted argument uses, and every step of the counted proof built on it is retained verbatim in the delivered body.\par\medskip
\ifdefined\hypertarget\hypertarget{sourcepointer-3}{}\fi
\noindent\textit{Source reference.} Original equation 08H7: sources/stacks-spaces-perfect/pinned-original-source.tex, source lines 320--320. Prerequisite pointer invoked at proof line(s) 3437; it names a standard fact of the same chapter that the counted argument uses, and every step of the counted proof built on it is retained verbatim in the delivered body.\par\medskip
\ifdefined\hypertarget\hypertarget{sourcepointer-4}{}\fi
\noindent\textit{Source reference.} Original reference 08H8: sources/stacks-spaces-perfect/pinned-original-source.tex, source lines 330--330. Prerequisite pointer invoked at proof line(s) 3441; it names a standard fact of the same chapter that the counted argument uses, and every step of the counted proof built on it is retained verbatim in the delivered body.\par\medskip
\ifdefined\hypertarget\hypertarget{sourcepointer-5}{}\fi
\noindent\textit{Source reference.} Original lemma 06Z0: sources/stacks-chow/tags.txt, source lines 8904--8904. Prerequisite pointer invoked at proof line(s) 3486; it names a standard fact of a sibling chapter that the counted argument uses, and every step of the counted proof built on it is retained verbatim in the delivered body.\par\medskip
\ifdefined\hypertarget\hypertarget{sourcepointer-6}{}\fi
\noindent\textit{Source reference.} Original lemma 08EE: sources/stacks-chow/tags.txt, source lines 10670--10670. Prerequisite pointer invoked at proof line(s) 3452; it names a standard fact of a sibling chapter that the counted argument uses, and every step of the counted proof built on it is retained verbatim in the delivered body.\par\medskip
\ifdefined\hypertarget\hypertarget{sourcepointer-7}{}\fi
\noindent\textit{Source reference.} Original lemma 08EB: sources/stacks-chow/tags.txt, source lines 10667--10667. Prerequisite pointer invoked at proof line(s) 3487; it names a standard fact of a sibling chapter that the counted argument uses, and every step of the counted proof built on it is retained verbatim in the delivered body.\par\medskip
\ifdefined\hypertarget\hypertarget{sourcepointer-8}{}\fi
\noindent\textit{Source reference.} Original lemma 08JI: sources/stacks-chow/tags.txt, source lines 10851--10851. Prerequisite pointer invoked at proof line(s) 3463; it names a standard fact of a sibling chapter that the counted argument uses, and every step of the counted proof built on it is retained verbatim in the delivered body.\par\medskip
\ifdefined\hypertarget\hypertarget{sourcepointer-9}{}\fi
\noindent\textit{Source reference.} Original remark 08JF: sources/stacks-chow/tags.txt, source lines 10848--10848. Prerequisite pointer invoked at proof line(s) 3445; it names a standard fact of a sibling chapter that the counted argument uses, and every step of the counted proof built on it is retained verbatim in the delivered body.\par\medskip
\begin{lemma}
\label{lemma-descend-RHom}
Let $X$ be a scheme. Let $E, F$ be objects of $D(\mathcal{O}_X)$.
Assume either
\begin{enumerate}
\item $E$ is pseudo-coherent and $F$ lies in $D^+(\mathcal{O}_X)$, or
\item $E$ is perfect and $F$ arbitrary,
\end{enumerate}
then there is a canonical isomorphism
$$
\epsilon^*R\SheafHom(E, F) \longrightarrow R\SheafHom(\epsilon^*E, \epsilon^*F)
$$
Here $\epsilon$ is as in (\ref{equation-epsilon}).
\end{lemma}

\begin{proof}
Recall that $\epsilon$ is flat (Lemma \ref{lemma-epsilon-flat}) and
hence $\epsilon^* = L\epsilon^*$. There is a canonical map
from left to right by
Cohomology on Sites, Remark
\ref{sites-cohomology-remark-prepare-fancy-base-change}.
To see this is an isomorphism we can work locally, i.e., we may
assume $X$ is an affine scheme.

\medskip\noindent
In case (1) we can represent $E$ by a bounded above complex
$\mathcal{E}^\bullet$ of finite free $\mathcal{O}_X$-modules, see
Derived Categories of Schemes, Lemma \ref{perfect-lemma-lift-pseudo-coherent}.
We may also represent $F$ by a bounded below complex $\mathcal{F}^\bullet$
of $\mathcal{O}_X$-modules. Applying
Cohomology, Lemma
\ref{cohomology-lemma-Rhom-complex-of-direct-summands-finite-free}
we see that $R\SheafHom(E, F)$ is represented by the complex with terms
$$
\bigoplus\nolimits_{n = - p + q}
\SheafHom_{\mathcal{O}_X}(\mathcal{E}^p, \mathcal{F}^q)
$$
Applying Cohomology on Sites, Lemma
\ref{sites-cohomology-lemma-Rhom-complex-of-direct-summands-finite-free}
we see that $R\SheafHom(\epsilon^*E, \epsilon^*F)$ is represented by the
complex with terms
$$
\bigoplus\nolimits_{n = - p + q}
\SheafHom_{\mathcal{O}_\etale}
(\epsilon^*\mathcal{E}^p, \epsilon^*\mathcal{F}^q)
$$
Thus the statement of the lemma boils down to the true fact
that the canonical map
$$
\epsilon^*\SheafHom_{\mathcal{O}_X}(\mathcal{E}, \mathcal{F})
\longrightarrow
\SheafHom_{\mathcal{O}_\etale}
(\epsilon^*\mathcal{E}, \epsilon^*\mathcal{F})
$$
is an isomorphism for any $\mathcal{O}_X$-module $\mathcal{F}$ and
finite free $\mathcal{O}_X$-module $\mathcal{E}$.

\medskip\noindent
In case (2) we can represent $E$ by a strictly perfect
complex $\mathcal{E}^\bullet$ of $\mathcal{O}_X$-modules, use
Derived Categories of Schemes, Lemmas
\ref{perfect-lemma-affine-compare-bounded} and
\ref{perfect-lemma-perfect-affine} and the fact that a perfect
complex of modules is represented by a finite complex of finite
projective modules. Thus we can do the exact same proof as
above, replacing the reference to
Cohomology, Lemma
\ref{cohomology-lemma-Rhom-complex-of-direct-summands-finite-free}
by a reference to
Cohomology, Lemma
\ref{cohomology-lemma-Rhom-strictly-perfect}.
\end{proof}

\end{document}
